\documentclass[11pt]{article}
\usepackage{amsmath,amssymb,amsthm,hyperref}
\newtheorem{theorem}{Theorem}
\newtheorem{lemma}[theorem]{Lemma}
\newtheorem{corollary}[theorem]{Corollary}
\title{A Low-Degree Filtration for Boolean Reciprocal Lattice Jets}
\author{}
\date{November 23, 2026}
\begin{document}
\maketitle

\paragraph{Scope.}
This note gives a necessary condition for higher-order perturbations
of the reciprocal lattice used by repeated P\'olya averaging. It is
a manual proof, not a Lean formalization, and does not prove the
universal head bound. In particular, it does not discard higher
orders merely because first variations fail.

\section{The common reciprocal polynomial}

Let $n\ge1$, $k\ge0$, and $\delta>0$. Let $C(x)$ be affine and
strictly positive on $\{0,1\}^n$. For any nonnegative integer vector
$\alpha=(\alpha_0,\ldots,\alpha_n)$ with $|\alpha|\le k$, put
\[
 L_\alpha(x)=\alpha_0+\sum_i\alpha_i x_i,
 \qquad B_\alpha(x)=C(x)+\delta L_\alpha(x),
 \qquad Q(t)=\prod_{j=0}^k(t+\delta j).
 \tag{1}
\]
All these denominators are positive on the cube. Requiring the
additional same-sign slope condition only restricts the family.
Let $\mathcal P_d$ be the space of Boolean tables represented by
real polynomials of degree at most $d$, and set
\[
            D_d=\dim\mathcal P_d
                  =\sum_{j=0}^{\min(d,n)}\binom nj.
 \tag{2}
\]
Every polynomial table has a unique multilinear representative;
reducing $x_i^m$ to $x_i$ never increases degree.

\begin{lemma}[Division and differentiation]
There is a polynomial $R(t,z)$ of total degree at most $k$ such that
$R(t,j)=Q(t)/(t+\delta j)$ for $j=0,\ldots,k$. Define
$W_\ell(t)=(d/dt)^\ell(1/Q(t))$. For every integer $m\ge0$,
every eligible $\alpha$, and every Boolean input,
\[
 \frac1{B_\alpha(x)^{m+1}}
 =\frac{(-1)^m}{m!}
   \sum_{j=0}^{\min(m,k)}\binom mj
       (\partial_t^j R)(C(x),L_\alpha(x))W_{m-j}(C(x)).
 \tag{3}
\]
\end{lemma}
\begin{proof}
Divide $Q(t)$ by the monic $t+\delta z$ in $\mathbb R[z][t]$.
Its quotient $R$ has total degree at most $k$. At $z=j$ the
remainder vanishes, proving the first statement. Fix a Boolean
input. Its value $L_\alpha(x)$ is an integer between zero and $k$,
so the corresponding identity of rational functions in $t$ is
valid before substituting $t=C(x)$. Differentiate it $m$ times
and use the product rule. The $j$th derivative of $R$ has degree
at most $k-j$ and is zero for $j>k$, giving (3).
\end{proof}

\section{All coefficients up to a specified perturbation order}

For each of finitely many heads choose a base $B_{\alpha_h}$ from
(1). Let $A_h(\varepsilon,x)$ and
$B_h(\varepsilon,x)$ be power series in $\varepsilon$ whose
coefficients are affine functions of $x$, with
$B_h(0,x)=B_{\alpha_h}(x)$. Also allow a scalar power-series bias.
These can be formal series for the coefficient assertion, or
convergent real analytic series for the sign assertion. Write
\[
 F(\varepsilon,x)=c(\varepsilon)+
                       \sum_h A_h(\varepsilon,x)/B_h(\varepsilon,x).
 \tag{4}
\]
No cancellation identity, equal root multiset, or restriction on the
numerator coefficients is assumed.

\begin{theorem}[Jet filtration]
Every coefficient of (4) through order $r$ belongs to
\[
 \mathcal S_{k,r}(C,\delta)
       =\sum_{\ell=0}^r W_\ell(C(x))\mathcal P_{k+\ell+1}.
 \tag{5}
\]
In particular,
\[
 \dim\mathcal S_{k,r}\le\sum_{\ell=0}^r D_{k+\ell+1}
                   \le(r+1)D_{k+r+1}.
 \tag{6}
\]
If each input has a nonzero coefficient at some order at most $r$,
the strict sign table of $F(\varepsilon,\cdot)$ for all sufficiently
small positive $\varepsilon$ is represented by an element of (5).
The first nonzero orders may differ between inputs.
\end{theorem}
\begin{proof}
Expand a reciprocal denominator formally by its geometric series
about $B_{\alpha_h}$. A coefficient at order at most $r$ is a
linear combination of terms
\[
                    P_m(x)/B_{\alpha_h}(x)^{m+1},
             \qquad 0\le m\le r,\quad \deg P_m\le m+1.
 \tag{7}
\]
There are $m$ affine denominator-increment factors and at most one
affine numerator factor, proving the degree bound. In the $j$th
summand of (3), multiplication by $P_m$ gives a polynomial factor
of degree at most $m+1+k-j=k+(m-j)+1$. Setting $\ell=m-j$
puts this term in the $\ell$th summand of (5). A scalar bias
belongs to its $\ell=0$ summand, since $1=Q(C)/Q(C)$ and
$\deg Q(C)\le k+1$. This proves the coefficient assertion and (6).

For convergent series, at each input the sign for small positive
$\varepsilon$ is the sign of its first nonzero coefficient. The
same holds for the degree-$r$ Taylor truncation. There are finitely
many inputs, so one sufficiently small positive number works for
all of them. The truncation at that number is a linear combination
of the first $r+1$ coefficient tables, hence lies in (5).
\end{proof}

For constant-numerator first variations one obtains the sharper
space bound
\[
 \dim\left\{c+\sum_\alpha D_\alpha(x)/B_\alpha(x)^2:
                         D_\alpha\text{ affine}\right\}
                    \le 1+D_k+D_{k+1}.
 \tag{8}
\]
Indeed (3) with $m=1$ gives a sum of
$W_0(C)\mathcal P_k$ and $W_1(C)\mathcal P_{k+1}$, plus the bias.
For $k=1$, this is $\binom n2+2n+3$, agreeing with the direct
one-step calculation in \texttt{polya-reciprocal-audit.tex}.

\section{Counting target-dependent base parameters}

\begin{theorem}[Low-order lattice jets are not universal]
Fix $\eta>0$. Suppose $k,r\ge0$ and
$k+r+1\le(1/2-\eta)n$. Even allowing $C$, $\delta$, all selected
lattice bases, all numerator series, and all denominator series to
depend on the target, the stable strict sign tables in the preceding
theorem number only $2^{o(2^n)}$ as $n$ tends to infinity.
The conclusion is uniform over these integers $k,r$ and does not
require a bound on the number of selected lattice bases.
\end{theorem}
\begin{proof}
Write $W_\ell(t)=N_\ell(t)/Q(t)^{\ell+1}$, with $N_0=1$.
Differentiation gives
\[
 N_{\ell+1}=Q N_\ell'-(\ell+1)Q'N_\ell.
 \tag{9}
\]
Regarding $t,\delta$ as variables, $Q$ is homogeneous of degree
$k+1$ and induction gives total degree at most $\ell k$ for
$N_\ell$. For an arbitrary table in (5), clear the positive
denominator $Q(C)^{r+1}$. Each resulting input polynomial has
degree at most one in the
\[
                       p_1=\sum_{\ell=0}^r D_{k+\ell+1}
\]
free coefficients of the multilinear polynomial factors in (5).
Its degree in the $q=n+2$ coefficients of $C$ and the step
$\delta$ is at most $r(k+1)$: the factor multiplying the
$\ell$th polynomial is $N_\ell(C)Q(C)^{r-\ell}$, whose degree
is at most $\ell k+(r-\ell)(k+1)$.

Put $p=p_1+q$. Once $p\le N=2^n$, the two-block sign bound from
\texttt{mechanism-audit.tex} gives
\[
 \left(\frac{32768N}{p}\right)^p
      2^{p_1}\max\{2,r(k+1)\}^{q}.
 \tag{10}
\]
This is an enlarged family: arbitrary polynomial coefficients in
(5) need not be realizable by actual paths, but every actual
coefficient or stable sign table under consideration is included.

If $a=1/2-\eta$ is nonnegative, the elementary binomial-entropy
bound gives $D_{\lfloor an\rfloor}\le2^{h_2(a)n}$, where
$h_2(a)=-a\log_2a-(1-a)\log_2(1-a)<1$ (with the usual endpoint
convention). Thus $p_1\le(n+1)2^{h_2(a)n}$, and every logarithm
in (10) is $o(N)$. For $a<0$ the hypotheses have no instances.
There are at most $(n+1)^2$ possible $(k,r)$, so their union still
has $2^{o(N)}$ sign tables. No factor for choosing the individual
$\alpha_h$ is needed: every choice is already included in the
single envelope (5).
\end{proof}

\begin{corollary}[A necessary linear perturbation order]
Let $\alpha_*>0$ be the unique solution of
\[
       (1+\alpha_*)\log_2(1+\alpha_*)
                         -\alpha_*\log_2\alpha_*=1.
 \tag{11}
\]
It is approximately $0.29381537$. Consider complete depth-$k$
layers, charging all $\binom{n+k}{k}$ distinct terminal forms.
If this cost is $O(2^n/n^2)$, a universal construction by the
analytic paths above cannot have maximum first-nonzero order at
most $\beta n$ for any fixed $\beta<1/2-\alpha_*$.
The numerical value of this threshold is approximately $0.20618463$.
\end{corollary}
\begin{proof}
The cost hypothesis and the factorial estimate imply
$k/n\le\alpha_*+o(1)$. Indeed the exponential rate of
$\binom{n+k}{k}$, when $k/n$ tends to a positive constant
$\alpha$, is the strictly increasing expression on the left of
(11). A larger limiting ratio would exceed rate one; unbounded
$k/n$ also violates the cost bound. If $\beta<1/2-\alpha_*$,
choose a fixed positive gap $\eta$ smaller than
$1/2-\alpha_*-\beta$. For sufficiently large $n$ the preceding
theorem applies, excluding universality.
\end{proof}

This is a necessary order bound, not a construction at or above
that order. The polynomial coefficients in (5) are only an upper
envelope, and their compatibility with one collection of affine
head paths has not been proved. Assigning them independently and
calling the resulting table a head score would be an invalid step.
Arbitrary higher orders and general admissible base configurations
remain outside this obstruction.

\end{document}
